% ICAIF 2026 Submission - ACM sigconf format
% Use: \documentclass[sigconf,anonymous]{acmart} for submission
% Use: \documentclass[sigconf]{acmart} for camera-ready

\documentclass[sigconf,anonymous]{acmart}

% Remove copyright for submission
\setcopyright{none}
\settopmatter{printacmref=true, printfolios=false}
\renewcommand\footnotetextcopyrightpermission[1]{}
\pagestyle{plain}
\acmConference[]{}{}{}

% Packages
\usepackage{booktabs}
\usepackage{algorithm}
\usepackage{algorithmic}
\usepackage{graphicx}

% BibTeX
\bibliographystyle{ACM-Reference-Format}

\begin{document}
\sloppy
\emergencystretch 1.5em

\title{Regime-Dependent Predictive Structure Between Equity Factors: Evidence from Granger Causality}

\author{Chorok Lee}
\authornote{Corresponding author}
\email{chorok.lee@kaist.ac.kr}
\affiliation{%
  \institution{Korea Advanced Institute of Science and Technology (KAIST)}
  \city{Daejeon}
  \country{South Korea}
}

\begin{abstract}
Do predictive relationships between equity factors change across market regimes?
Using 35 years of daily Fama-French data (1990--2024) and a Student-$t$ Hidden
Markov Model, we find that HML$\to$SMB Granger causality is regime-dependent and
era-shifting.
In-sample, the link is Bonferroni-significant in the Normal regime pre-GFC
($p = 8.75 \times 10^{-9}$, $\Delta R^2 = 2.06\%$) but vanishes after a
structural break at January 2008 (Chow $p = 2.29\times10^{-6}$; post-2008:
$p = 0.73$, TOST-confirmed equivalent to zero).
A frozen out-of-sample design (HMM trained 1990--2012, tested 2013--2024) detects
modest Elevated-regime evidence (permutation $p = 0.022$ [primary],
$0.063$ [sensitivity bound]; 2 of 3 HMM optima significant), but this does not
survive 3-regime Bonferroni correction and requires $K = 3$.
Economic magnitudes are small ($\Delta R^2 < 1\%$) and do not translate to
trading profitability (Sharpe $= -0.07$).
Transfer entropy reveals a reverse nonlinear SMB$\to$HML channel ($z = 5.37$),
indicating the predictive architecture is directionally asymmetric: linear forward,
nonlinear reverse.
Secondary validation via MOM$\to$SMB (permutation $p = 0.010$) confirms that
regime-conditional predictive structure extends beyond a single factor pair.
We interpret all findings as predictive precedence, not structural causality.
\end{abstract}

\begin{CCSXML}
<ccs2012>
   <concept>
       <concept_id>10002950.10003648.10003671</concept_id>
       <concept_desc>Mathematics of computing~Time series analysis</concept_desc>
       <concept_significance>500</concept_significance>
   </concept>
   <concept>
       <concept_id>10010147.10010257.10010293.10010294</concept_id>
       <concept_desc>Computing methodologies~Causal reasoning and diagnostics</concept_desc>
       <concept_significance>500</concept_significance>
   </concept>
</ccs2012>
\end{CCSXML}

\ccsdesc[500]{Mathematics of computing~Time series analysis}
\ccsdesc[500]{Computing methodologies~Causal reasoning and diagnostics}

\keywords{Factor Investing, Granger Causality, Regime Switching, Hidden Markov Models, Model Complexity Characterization, Out-of-Sample Validation, Risk Management}

\maketitle

\section{Introduction}

The August 2007 quantitative meltdown, in which systematic equity strategies lost
more than 27\% in three trading days (August 7--9, 2007)~\cite{khandani2011quants}, revealed a critical
blind spot in factor-based risk management.
When multiple quantitative funds held similar factor exposures, forced liquidation
by one fund created price pressure that cascaded to all others.
Standard correlation-based risk models failed to anticipate this cascade because they
measure \emph{co-movement} but not \emph{temporal precedence}.

\textbf{The Missing Piece: Which Factor Moves First?}

Existing research establishes three stylized facts:
(1) Equity return correlations increase during market stress~\cite{ang2002asymmetric};
(2) Returns exhibit regime-switching behavior~\cite{hamilton1989new};
(3) Factor crowding amplifies drawdowns during liquidation~\cite{stein2009presidential}.

However, a critical question remains underexplored: \textbf{Does the predictive
architecture between factors change qualitatively across market regimes?}

Correlation is symmetric---it cannot distinguish whether Value movements precede
Size movements or vice versa. Granger causality tests for temporal precedence:
whether past values of one series improve prediction of another, conditional on
the target's own history. If such predictive relationships vary by regime, we could:
\begin{itemize}
    \item Monitor the leading factor to anticipate movements in the lagging factor
    \item Adjust hedges based on the current regime's predictive structure
    \item Detect regime transitions by observing when predictive links emerge or disappear
\end{itemize}

\textbf{Important caveat.} Granger causality establishes \emph{predictive precedence},
not \emph{structural causality}. A statistically significant Granger relationship
could arise from: (1) direct causal influence, (2) a common driver affecting both
series with different lags, or (3) omitted variables. We interpret our findings as documenting regime-dependent predictive structure,
with economic mechanisms offered as hypotheses rather than verified causal channels---a
distinction foundational in the causal inference literature~\cite{pearl2009causality}
and articulated for factor research by L\'opez de Prado~\cite{lopezdeprado2023causal}.
We adopt the Granger framework for its tractability on high-dimensional daily
factor series. L\'opez de Prado advocates structural causal models (SCMs/DAGs)
as the more rigorous alternative for factor research; we acknowledge this
distinction throughout.

\textbf{Hidden Markov Model framework.} We use a Hidden Markov Model---a
probabilistic framework that identifies latent market regimes based on return
distributions---to condition Granger tests on market state. Unlike ad hoc crisis
windows or fixed breakpoints, the HMM discovers regimes unsupervised from
distributional properties alone, reducing researcher degrees of freedom.

\subsection{Our Findings}

Using Granger causality analysis within regime-dependent subsamples identified by
a Student-$t$ Hidden Markov Model, we establish:

\textbf{Central Finding: Regime-dependent, era-shifting predictive structure.}
In-sample analysis (1990--2024) reveals that HML$\to$SMB is Bonferroni-significant in
the Normal regime pre-GFC ($p = 8.75 \times 10^{-9}$, HAC $p = 9.45 \times 10^{-8}$,
$\Delta R^2 = 2.06\%$) but undergoes a structural break at January 2008
(Chow $p = 2.29\times10^{-6}$; Bai-Perron sup-$F = 22.1$; post-2008: $p = 0.73$).
In-sample Elevated-regime HML$\to$SMB is suggestive ($p = 0.004$, below Bonferroni
threshold), and a single-stage Markov-Switching VAR further supports the historical
pattern ($p < 10^{-12}$; Appendix~\ref{app:msvar}). This is consistent with the
hypothesis that predictive content migrated from Normal to Elevated volatility
states post-GFC, though regime redistribution under the frozen classifier is a
plausible alternative explanation (Section~\ref{sec:frozen_oos}).

\textbf{OOS evidence in the Elevated regime.}
Using a frozen OOS design (HMM trained on 1990--2012, Granger tested on
2013--2024), HML$\to$SMB shows modest evidence in the Elevated regime---the
same regime that shows in-sample suggestive significance. The OOS test has
98.7\% power to detect the in-sample effect size, so the smaller OOS magnitude
($\Delta R^2 = 0.73\%$ vs.\ $2.06\%$) reflects genuine attenuation, not low power.
This is consistent with the migration hypothesis, though regime redistribution
under the frozen classifier (Elevated grows from 13.7\% to 30.7\% of OOS days)
is a plausible confound
(economically motivated pair; ranks 2nd of 30 directed pairs by $F$-statistic;
see Section~\ref{sec:pairselection}):
\begin{itemize}
    \item \emph{Elevated (OOS):} $F$-$p = 0.014$, HAC $p = 0.041$ (Newey--West, bandwidth~1),
    $\Delta R^2 = 0.73\%$ (lag 1, pre-fixed from in-sample BIC); permutation
    $p = 0.022$ (primary; model-free, 50{,}000 shuffles~\cite{good2005permutation}),
    $p = 0.063$ (sensitivity bound; decimal-unit, 10{,}000 shuffles);
    significant in 2 of 3 distinct local HMM optima.
    \item \emph{Normal (OOS):} $p = 0.149$ (null).
    \item \emph{Crisis (OOS):} $p = 0.281$ (null).
\end{itemize}

\textbf{Forward direction predominantly linear; reverse direction nonlinear.}
A four-model diagnostic (Linear, RF, MLP, LSTM; 200 permutations) finds
no nonlinear method improves forward HML$\to$SMB prediction beyond the linear
baseline under the primary fit. However, transfer entropy detects a
reverse nonlinear SMB$\to$HML channel in the Normal regime ($z = 5.37$,
$p < 10^{-6}$; Section~\ref{sec:robustness}), suggesting the overall
predictive architecture is mixed---linear forward, nonlinear reverse.
The ``predominantly linear'' characterization applies to the forward direction
under the primary HMM solution; nonlinear structure emerges under an
alternative solution (Appendix~\ref{app:robustness}).

\subsection{Contributions}

\begin{enumerate}
    \item \textbf{Regime-dependent predictive structure with a dated structural break.}
    We document that HML$\to$SMB Granger causality is not merely time-varying
    but \emph{regime-specific}: Bonferroni-significant in the Normal regime
    pre-GFC ($p = 8.75 \times 10^{-9}$, $\Delta R^2 = 2.06\%$ above the SMB
    AR baseline), absent post-2008 (Chow $p = 2.29\times10^{-6}$; TOST
    $p_{\text{equiv}} < 0.001$). This challenges the implicit assumption in
    factor-timing models that cross-factor predictive links are
    regime-invariant (Section~\ref{sec:main_result}).

    \item \textbf{Modest OOS evidence with disclosed fragility.}
    A frozen OOS design detects Elevated-regime HML$\to$SMB
    (permutation $p = 0.022$ [primary], $0.063$ [sensitivity bound];
    Table~\ref{tab:frozen_granger}). We designate permutation as primary
    (model-free, no distributional assumptions).
    The result requires $K = 3$, does not survive 3-regime Bonferroni,
    and regime redistribution is a plausible confound.
    A secondary pair (MOM$\to$SMB, permutation $p = 0.010$) independently
    confirms the regime-conditional pattern
    (Sections~\ref{sec:frozen_oos}, \ref{sec:pairselection}, \ref{sec:robustness}).

    \item \textbf{Directional complexity characterization.}
    A four-model diagnostic (OLS, RF, MLP, LSTM) finds no nonlinear
    improvement for forward HML$\to$SMB under the primary fit, but transfer
    entropy detects a reverse nonlinear SMB$\to$HML channel ($z = 5.37$
    in Normal). This directional asymmetry---linear forward, nonlinear
    reverse---suggests the underlying information flow is richer than
    either Granger or TE alone captures
    (Sections~\ref{sec:hmm}, \ref{sec:robustness}).
    A hybrid VaR detector (Appendix~\ref{app:var}) demonstrates that
    regime-conditional architectures can improve tail-risk monitoring,
    though the improvement is driven by the volatility-override mechanism
    rather than the Granger signal itself.
\end{enumerate}

\section{Related Work}

Our work intersects four literatures. We summarize each and identify the gap we address.

\textbf{Factor time-series predictability.}
Fama and French~\cite{fama1993common,fama2015five} establish that Size, Value,
Profitability, and Investment factors explain cross-sectional return variation.
Ehsani and Linnainmaa~\cite{ehsani2022factor} and Gupta and Kelly~\cite{gupta2019factor}
show that factor returns exhibit time-series momentum (within-factor autocorrelation);
Haddad, Kozak and Santosh~\cite{haddad2020factor} demonstrate factor-level predictability
from SDF-motivated conditioning variables.
Chordia and Shivakumar~\cite{chordia2002momentum} and Cooper et al.~\cite{cooper2004market}
establish regime-conditional momentum using observable macro states.
These papers examine within-factor momentum but not \emph{cross-factor Granger structure}---which
factor leads the other, under which regime. Our Granger test conditions on SMB's own lags,
isolating incremental cross-factor content invisible to unconditional approaches
(see also~\cite{jegadeesh1993returns,carhart1997persistence,barroso2015momentum,daniel2016momentum,flogel2022momentum,neuhierl2024timing,asness2013value}).

\textbf{Factor crowding and systemic risk.}
Shleifer and Vishny~\cite{shleifer1997limits} establish that capital constraints create
return spillovers; Marks and Shang~\cite{marks2019factor} and Chincarini et al.~\cite{chincarini2026crowded}
provide evidence that HML/SMB crowding leads to crash risk.
Betermier et al.~\cite{betermier2025investor} document institutional portfolio tilts
generating common HML/SMB exposures; Brunnermeier and Pedersen~\cite{brunnermeier2009market}
formalize margin-constraint spirals; Coval and Stafford~\cite{coval2007asset} document
liquidation-driven fire sales
(see also~\cite{stein2009presidential,anton2014connected,lou2022comomentum,kang2021crowding,volpati2020zooming,hua2024dynamics,greenwood2011stock}).
Existing work measures crowding \emph{intensity} but not
\emph{regime-conditional predictive spillover} between factors.

\textbf{Regime-switching models.}
Hamilton~\cite{hamilton1989new} introduced Markov-switching models;
Krolzig~\cite{krolzig1997markov} extends to MS-VARs (the single-stage analog we
compare against in Section~\ref{sec:msvar}).
Bulla~\cite{bulla2011hidden} introduces Student-$t$ HMMs for financial returns;
Chincoli and Guidolin~\cite{chincoli2017linear} show MS-VARs improve factor-return
predictability without examining causal direction.
Shu and Mulvey~\cite{shu2025dynamic} and Nuriyev et al.~\cite{nuriyev2024regime}
apply regime signals to factor \emph{allocation}---which factor to hold per regime---whereas
we examine Granger \emph{precedence between} factors within shared regimes
(see also~\cite{ang2002regime,guidolin2007asset,guidolin2008size,oelschlager2024fhmm,ang2012regime,baitinger2024simplicity,shu2024downside,wang2020regime}).

\textbf{Granger causality, connectedness, and complexity.}
Lo and MacKinlay~\cite{lo1990contrarian} establish cross-autocorrelation
between large-cap and small-cap portfolios; Granger~\cite{granger1969investigating}
introduced the predictive-precedence test we employ.
Psaradakis et al.~\cite{psaradakis2005markov} pioneer regime-switching Granger causality;
Shi et al.~\cite{shi2018change,shi2020causal} detect causal change-points with recursive tests;
Droumaguet et al.~\cite{droumaguet2017granger} extend to Bayesian MS-VAR inference.
Billio et al.~\cite{billio2012econometric} construct Granger networks among financial
institutions; Diebold and Yilmaz~\cite{diebold2012better} develop FEVD spillover measures
(Buchwalter et al.~\cite{buchwalter2025clustered} show within-cluster FEVD limitations).
Tank et al.~\cite{tank2022neural} extend Granger causality to nonlinear settings;
Howard et al.~\cite{howard2025causal} apply causal discovery at the stock level.
We apply neural Granger, Random Forest, and transfer
entropy~\cite{schreiber2000measuring,marschinski2002analysing,frenzel2007partial,barnett2009granger}
to characterize the complexity boundary of the HML$\to$SMB link.
No prior work examines \emph{regime-conditional} Granger causality
between Fama-French equity \emph{factors}.

\section{Methodology}

\subsection{Data}

We use daily returns for six equity factors from Kenneth French's data library~\cite{french2024data} (five Fama-French factors~\cite{fama2015five} plus Carhart Momentum):
Market excess return (MKT-RF), Size (SMB), Value (HML), Profitability (RMW),
Investment (CMA), and Momentum (MOM)~\cite{jegadeesh1993returns,carhart1997persistence}.

\textbf{Sample:} January 2, 1990 -- December 31, 2024 (8,817 trading days).
Factor returns are used as downloaded (daily percentages; e.g., 0.10 = 0.1\%).
Granger $F$-statistics are scale-invariant. The HMM fits and Granger tests for the
in-sample analysis use percentage units. The frozen OOS pipeline converts to decimal form
before HMM fitting and Granger testing (e.g., 0.001 $= 0.1\%$); the frozen HMM is trained
on rescaled data, so all computations within that pipeline are internally consistent.
This scale difference does not affect conclusions: Granger $F$-statistics and HAC $t$-ratios
are scale-invariant, and the OOS Elevated $F$-$p < 0.05$ under both conventions.
Daily factor returns are stationary by construction; ADF tests confirm unit-root
rejection ($p < 0.001$) for all six series (BIC lag selection, intercept included,
no trend).

\subsection{Student-\texorpdfstring{$t$}{t} Hidden Markov Model}
\label{sec:hmm}

Let $z_t \in \{1, \ldots, K\}$ denote the latent regime at time $t$. We model:

\textbf{Transition model:}
$P(z_t = k \mid z_{t-1} = j) = A_{jk}$

\textbf{Emission model (multivariate Student-$t$):}
\begin{equation}
p(\mathbf{x}_t \mid z_t = k) \propto \left(1 + \frac{\delta_k(\mathbf{x}_t)}{\nu_k}\right)^{-\frac{\nu_k+d}{2}}
\end{equation}

where $\delta_k(\mathbf{x}_t) = (\mathbf{x}_t - \boldsymbol{\mu}_k)^\top\boldsymbol{\Sigma}_k^{-1}(\mathbf{x}_t - \boldsymbol{\mu}_k)$.

\textbf{Why Student-$t$?} Financial returns exhibit excess kurtosis. Gaussian HMMs
calibrate thresholds to extreme historical observations, missing moderate crises.
Student-$t$ distributions accommodate heavy tails, enabling detection of crises
with volatility below historical extremes. Estimated degrees of freedom for the
primary fit (seed 28): $\hat{\nu}_{\text{Normal}} = 6.2$, $\hat{\nu}_{\text{Elevated}} = 3.9$,
$\hat{\nu}_{\text{Crisis}} = 5.5$---all consistent with heavy tails well below
the Gaussian limit. The Elevated regime ($\hat\nu = 3.9 < 4$) has technically
infinite kurtosis, reflecting extreme tail behavior in stress-adjacent periods;
the HAC-robust Granger tests remain valid since they require no finite-moments
assumption beyond what is needed for asymptotic normality of OLS coefficients.

\textbf{Model selection.} The number of regimes ($K=3$) was selected by the Bayesian
Information Criterion, which favored three regimes over two
($\Delta\text{BIC} = \text{BIC}(K{=}2) - \text{BIC}(K{=}3) = 847$) or
four ($\Delta\text{BIC} = \text{BIC}(K{=}4) - \text{BIC}(K{=}3) = 312$). Transition probabilities are freely estimated via the Expectation-Maximization algorithm.\footnote{We fit Student-$t$ HMMs with 50
EM initializations.
\emph{Primary selection (leakage-safe):} the fit with highest log-likelihood across
all 50 starts, regardless of event detection. Three seeds (28, 20, 6) reach the same
full-sample log-likelihood; seed 28 is selected as primary by convention
(first in sorted order). Note: when re-trained on 1990--2012 only (frozen OOS),
these seeds converge to distinct local optima with different train-period log-likelihoods.
This fit assigns 0\% of 2008, 2011, and
2020 stress windows to the Crisis regime (Table~\ref{tab:detection}).
\emph{Sensitivity:} among the 16/50 fits (32\%) that assign $\geq$50\% of 2008
days to Crisis, we separately report the highest-LL fit; the LL cost is $<$0.3\%.
All tables use the primary fit unless explicitly noted.
The FF 25 portfolio overlap analysis (Section~\ref{sec:overlap_evidence}) and VaR
monitoring (Section~\ref{sec:risk_monitoring}) use the sensitivity fit (seed 42),
which assigns 47.8\% of COVID-2020 days to Crisis and yields qualitatively identical
full-sample Granger results to the primary fit (see FF25 footnote and VaR section
for frozen OOS Granger disclosure); this is disclosed at each analysis.
Table~\ref{tab:detection} compares against a Gaussian HMM baseline
(also multi-start, 10 seeds).}

\textbf{Local optima and economic validity.}
Our 50-seed multistart reveals 7 distinct local optima clusters (Appendix~\ref{app:optima}).
We adopt a two-stage selection criterion: (1)~among fits with $\geq$50\% detection
of the 2008 GFC, (2)~select the highest log-likelihood. This trades statistical
parsimony for economic interpretability. The OOS Elevated results are robust
across all 7 clusters ($p < 0.05$ in 5/7).

\subsection{Sample Overlap Consideration}

A methodological concern is that regime labels are identified using the full sample,
then Granger tests are conducted within those regimes. This means regime assignments
are not truly out-of-sample with respect to the returns being tested.

We address this concern in three ways. First, the HMM uses \emph{distributional}
properties (mean, covariance, tail behavior) to assign regimes, while Granger
tests examine \emph{temporal dynamics} (lagged predictive relationships). These
are distinct statistical features, limiting circularity. Second, our event-based
validation (Section~\ref{sec:oos}) tests whether the in-sample pattern generalizes
to specific historical episodes, providing partial out-of-sample assessment.
Third, a soft-label sensitivity check weighting observations by posterior regime
probabilities (rather than hard Viterbi assignments) yields qualitatively
identical Granger conclusions, confirming that regime-boundary uncertainty
does not drive the results.

We adopt the fully rigorous approach as the \emph{primary} validation,
following the OOS evaluation standard of Welch and Goyal~\cite{welch2008comprehensive}:
the HMM is fitted on 1990--2012 and all parameters are frozen before
classifying 2013--2024; Granger tests are then run on the held-out labels
using lag $= 1$ pre-fixed from in-sample BIC
(Section~\ref{sec:frozen_oos}). This frozen OOS design eliminates circularity
and yields significant HML$\to$SMB in the Elevated regime of the test period.
Additionally, BIC model selection ($K=3$ vs.\ $K=2$ vs.\ $K=4$) is verified on
training data alone: $K=3$ achieves BIC~$= 44{,}652$ vs.\ $K=2$'s $46{,}332$
($\Delta\text{BIC} = 1{,}680$), confirming $K=3$ is BIC-optimal without
accessing OOS data. (Note: the $\Delta\text{BIC} = 847$ reported in
Section~\ref{sec:hmm} is from full-sample model selection on 1990--2024;
the $\Delta\text{BIC} = 1{,}680$ here is from training-only selection on
1990--2012. Both strongly favor $K=3$.)

\subsection{Per-Regime Granger Causality}
\label{sec:granger_method}

Granger causality~\cite{granger1969investigating,shojaie2022granger} tests whether past values
of one series improve prediction of another, conditional on the target's own history---a
\emph{predictive precedence} criterion, not structural causality.
For each regime $k$, we extract observations: $\mathcal{T}_k = \{t : \hat{z}_t = k\}$.
For each ordered pair of factors $(i, j)$, we test:\footnote{We test pairwise
(bivariate) Granger relationships. A high-dimensional VAR with all 6 factors and 9 lags
would require $6 \times 6 \times 9 = 324$ parameters per regime; with the smallest regime
at $n \approx 1{,}000$ observations this is severely under-identified. Post-double-selection
methods~\cite{hecq2023granger} could address this in future work.}
$H_0: r_{j,t} \perp \{r_{i,t-\ell}\}_{\ell=1}^{L} \mid \{r_{j,t-\ell}\}_{\ell=1}^{L}$

\textbf{Multiple testing framework.}
In-sample analysis (1990--2024) uses a stringent Bonferroni correction:
$\alpha_{\text{fam}} = 0.01$ across 30 directed factor pairs, yielding
$\alpha_{\text{ind}} = 0.00033$. The frozen OOS analysis (2013--2024) is treated
as a confirmatory test of the pre-identified HML--SMB pair (economic prior:
Section~\ref{sec:pairselection}); given the design's separation from in-sample
pair selection, we apply $\alpha_{\text{OOS}} = 0.05$ for this single pair.
Results not meeting their respective threshold are reported as exploratory. HAC standard errors use Newey--West
with bandwidth equal to the Granger lag (bandwidth~1) as the primary specification,
applied consistently across all Granger tests. As a cross-convention check,
under the percentage-unit specification ($n = 953$, residual autocorrelation near zero),
Andrews's~\cite{andrews1991heteroskedasticity} AR(1) plug-in bandwidth for the
OOS Elevated series is also~1, yielding HAC $p = 0.021$---consistent with
the primary Newey--West $p = 0.041$.
Within-convention bandwidth sensitivity (percentage-unit; bandwidths 1, 4, 7
(rule-of-thumb for $n \approx 953$: $\lfloor 0.75 n^{1/3} \rfloor = 7$), 10, and 15)
yields $p < 0.05$ throughout (range: 0.015--0.025), confirming robustness to bandwidth choice.

\textbf{Regime boundary handling.} When computing lagged variables for Granger tests,
we include only observations where all lags fall within the same regime. Specifically,
for an observation at time $t$ in regime $k$ with maximum lag $L$, we require
$\hat{z}_{t-\ell} = k$ for all $\ell \in \{1, \ldots, L\}$. The exclusion varies
by analysis: Table~\ref{tab:main} uses BIC-selected lags (lag 1 for all rows);
Table~\ref{tab:r2} uses 15 lags; neural Granger and transfer entropy analyses
(Appendix~\ref{app:robustness}) use 9 lags; the VaR application
(Appendix~\ref{app:var}) uses model-specific warm-up requirements.
All analyses use trading-day-indexed lags; higher-lag analyses yield lower per-table $n$
(e.g., Normal: $n = 4{,}496$ at lag~9 vs.\ $n = 4{,}697$ at lag~1).

\textbf{Quantified impact.} Lag-1 boundary exclusion is low: 0.67\% in-sample (59/8,817),
though higher at 7.4\% for frozen OOS (224/3,020) due to more frequent
regime transitions in 2013--2024. This conservative rule may underestimate
transition-day effects; full exclusion statistics are in Appendix~\ref{app:optima}.

\subsection{Factor Pair Selection}
\label{sec:pairselection}

We focus on the HML--SMB pair for three reasons. First, \emph{economic
plausibility}: Value and Size factors have documented overlap in institutional
holdings (Section~\ref{sec:interpretation}), providing an economically motivated
mechanism---when leveraged institutions must unwind value positions, the associated
size exposure is simultaneously deleveraged, creating a predictive lag from HML
to SMB. This economic hypothesis precedes and motivates our data analysis.
Second, \emph{in-sample screening}: a preliminary screen of all 30 directed
factor pairs found that HML--SMB shows the most distinctive regime-heterogeneous
pattern, consistent with the economic prior. Crucially, the primary (LL-only) fit
independently identifies HML--SMB as distinctive (Normal: $p = 1.0 \times 10^{-6}$,
Elevated: $p = 0.013$, Crisis: $p = 0.880$ at lag 5)---no 2008-based screening is
involved in pair selection.
Third, \emph{architectural contrast}: lag-1 evidence is strongest in Normal and
substantially weaker outside Normal under the Bonferroni-significant threshold.

\textbf{Pair selection and multiple comparisons.} The OOS frozen test targets
the HML--SMB pair based on an economic prior (liquidity-mediated contagion),
but was not formally pre-registered. We therefore report results under
multiple correction frameworks:
(i)~uncorrected ($\alpha = 0.05$): both $F$-$p = 0.014$ and HAC $p = 0.041$ are significant;
(ii)~3-regime Bonferroni ($\alpha/3 = 0.0167$): $F$-$p = 0.014$ survives, HAC $p = 0.041$ does not;
(iii)~30-pair Bonferroni ($\alpha/30 = 0.00167$): neither survives.
The Elevated regime was identified as the post-GFC locus from in-sample temporal
analysis prior to OOS evaluation, motivating interpretation~(i); however, because
this identification was itself data-driven, interpretation~(ii) may be more
appropriate. We present the OOS Elevated result as modest evidence whose
strength depends on the reader's prior about pair and regime pre-specification.
\emph{Pre-registration status:} the HML--SMB pair was not formally pre-registered;
this analysis is therefore exploratory with respect to pair selection, even
though the economic prior (liquidity-mediated contagion) preceded data analysis.
A pre-registered replication on independent data (e.g., international factor
returns) would provide definitive confirmation.
\textbf{Permutation test scope.} A 50{,}000-shuffle permutation test
assesses whether the OOS Elevated Granger $F$-statistic is extreme under the null
of no predictive content. Specifically, we randomly permute HML residual labels
within the Elevated regime (preserving regime counts and temporal structure of
SMB), recompute the Granger $F$-statistic under each permutation, and compute
$p$ as the fraction of permuted $F$-values exceeding the observed
$F$~\cite{good2005permutation}. This nonparametric test avoids distributional
assumptions about the $F$-statistic under fat-tailed HMM residuals.
Result: $p = 0.022$ (SE $= 0.0007$, 95\% CI [0.021, 0.023], $n = 953$ clean Elevated observations)
under the percentage-unit convention.
Under a decimal-unit implementation (same seed and lag, stricter boundary/scale handling),
$n = 836$ valid OOS observations and 10{,}000 permutations yield $p = 0.063$
(MC SE $\approx 0.0024$), which does not reach 5\%.
The permutation evidence thus ranges from $p = 0.022$ to $p = 0.063$ across scale
conventions, straddling the conventional significance threshold.
We report the percentage-unit result as primary because the frozen HMM was trained
on percentage-unit data, but treat the decimal-unit result as a robustness bound
that qualifies the strength of the permutation evidence.
\textbf{30-pair multiplicity.} To directly assess selection bias, we run Granger
tests for all 30 directed factor pairs in the OOS Elevated window (percentage-unit
convention, seed~28). HML$\to$SMB ranks \textbf{2nd of 30} by $F$-statistic
($F = 9.06$, $p = 0.003$); the rank-based max-statistic $p = 2/30 = 0.067$.
The top-ranked pair is MOM$\to$SMB ($F = 20.3$), which was not the
economically motivated target of this study.
The rank-2 position provides direct evidence against post-hoc selection of HML$\to$SMB
from a broad search over 30 pairs.

\section{Results}

\subsection{Regime Characteristics}

Table~\ref{tab:regimes} summarizes the three identified regimes. Figure~\ref{fig:timeline}
shows regime assignments over the sample period with major crisis events marked.

\begin{table}[t]
\centering
\caption[Regime statistics]{Regime Summary Statistics (1990--2024, full-sample primary fit). Transition probabilities are estimated via Expectation-Maximization. The frozen OOS analysis (Section~\ref{sec:frozen_oos}) uses a separate fit trained on 1990--2012 only, yielding different regime counts. Elevated and Crisis regimes have similar mean norms (2.03 vs.\ 2.09) but are distinguished by tail shape ($\nu = 3.9$ vs.\ $5.5$) and covariance structure.}
\label{tab:regimes}
\small
\setlength{\tabcolsep}{4pt}
\begin{tabular}{lccccc}
\toprule
Regime & Days & Prop. & Mean $\|\mathbf{x}\|$ (daily \%) & $\nu$ & $P(z_t=z_{t-1})$ \\
\midrule
Normal & 4,723 & 53.6\% & 0.98 & 6.2 & 0.994 \\
Elevated & 3,023 & 34.3\% & 2.03 & 3.9 & 0.991 \\
Crisis & 1,071 & 12.1\% & 2.09 & 5.5 & 0.993 \\
\bottomrule
\end{tabular}
\setlength{\tabcolsep}{6pt}
\end{table}

\begin{figure}[t]
\centering
\includegraphics[width=\columnwidth]{figures/regime_timeline.pdf}
\caption{Regime assignments over sample period (1990--2024). Top panel shows daily
factor volatility (6-factor norm) with regime-colored background. Bottom panel shows
regime classifications. Vertical dashed lines mark major stress events (LTCM/Russia 1998, dot-com 2000--02,
GFC 2008, COVID-2020, 2022 rate shock). Under the
primary (leakage-safe) fit, Crisis regime (red) reflects a high-kurtosis statistical
state; the 2008-screened sensitivity fit aligns it with calendar crises (see
Section~\ref{sec:gaussian_comparison}). Elevated regime (yellow) captures moderate
volatility periods.}
\Description{Two-panel timeline showing daily factor-volatility norm with background regime colors and event markers, and a lower panel with daily regime assignments from 1990 to 2024.}
\label{fig:timeline}
\end{figure}

\subsection{Gaussian vs. Student-\texorpdfstring{$t$}{t} Comparison}
\label{sec:gaussian_comparison}

\begin{table}[t]
\centering
\caption{Crisis Detection Comparison. Detection rate = proportion of event days
assigned to Crisis regime.}
\label{tab:detection}
\small
\begin{tabular}{lccc}
\toprule
Event & Period & Student-$t$ & Gaussian \\
\midrule
2008 Financial & Jul '08 -- Jun '09 & 0.0\% & 83.3\% \\
2011 EU Debt & Jul -- Oct 2011 & 0.0\% & 25.9\% \\
2020 COVID-19 & Feb -- Jun 2020 & 0.0\% & 81.7\% \\
\bottomrule
\end{tabular}
\end{table}

Under the leakage-safe primary selection rule, the Student-$t$ fit does not
label these benchmark windows as Crisis (0.0\%), while the Gaussian baseline
assigns substantial Crisis mass to 2008 and 2020. We therefore treat event-level
``detection'' as sensitivity-dependent and avoid confirmatory claims based solely
on this table.
\textbf{Terminology note.} Under the leakage-safe primary fit, the third regime
captures a distinct statistical state (high kurtosis, shifted covariance) rather
than exclusively named crisis events---Table~\ref{tab:detection} shows 0\% alignment
with 2008, 2011, and 2020 stress windows. We retain the label ``Crisis'' for
consistency with the HMM literature~\cite{bulla2011hidden}, but readers should
interpret Crisis-regime results as pertaining to this statistical state. The
2008-screened sensitivity fit does align with calendar crises; we report it
separately where noted.

\subsection{Main Result: Regime-Dependent Predictive Structure}
\label{sec:main_result}

\begin{table}[t]
\centering
\caption{Granger Causality Between HML and SMB by Regime. Significance threshold:
$p < 0.00033$ (Bonferroni-corrected for 30 factor pairs).}
\label{tab:main}
\small
\begin{tabular}{llccc}
\toprule
Regime & Direction & $p$-value & Lag & Significant \\
\midrule
Normal & HML $\to$ SMB & $\mathbf{8.75 \times 10^{-9}}$ & 1 & \textbf{Yes} \\
Normal & SMB $\to$ HML & $0.864$ & 1 & No \\
\midrule
Elevated & HML $\to$ SMB & $0.004$ & 1 & Suggestive$^*$ \\
Elevated & SMB $\to$ HML & $0.897$ & 1 & No \\
\midrule
Crisis & HML $\to$ SMB & $0.695$ & 1 & No \\
Crisis & SMB $\to$ HML & $0.294$ & 1 & No \\
\bottomrule
\multicolumn{5}{l}{\footnotesize $^*$Below 0.01 but above Bonferroni threshold $\alpha/30 = 0.00033$.}
\end{tabular}
\end{table}

\textbf{Key finding under Bonferroni-significant threshold:} only Normal-regime
HML$\to$SMB survives Bonferroni correction.
\begin{itemize}
    \item \textbf{Normal (Bonferroni-significant):} Strong lag-1 effect
    ($p = 8.75 \times 10^{-9}$), robust to HAC correction. The Granger test conditions
    on lagged SMB, so the HML coefficient reflects \emph{incremental cross-factor}
    predictive content beyond SMB's own history---not a within-SMB autocorrelation artifact.
    \item \textbf{Elevated (suggestive):} HML$\to$SMB is directional at
    conventional levels ($p = 0.004$) but does not pass the paper's
    multiple-testing threshold.
    \item \textbf{Crisis (null):} No detectable lag-1 Granger effect in either
    direction.
\end{itemize}

\textbf{Robustness to heteroskedasticity.}
Financial returns exhibit volatility clustering, which may inflate standard Granger
$F$-test significance. We recompute at lag 1 (BIC-optimal) using Newey--West~\cite{newey1987simple}
heteroskedasticity-and-autocorrelation-consistent (HAC) standard errors.
Normal-regime HML$\to$SMB survives HAC correction
($p = 9.45 \times 10^{-8}$ HAC vs.\ $p = 8.75 \times 10^{-9}$ standard).
Elevated HML$\to$SMB weakens under HAC (1990--2024 in-sample: $p = 0.034$), and Crisis remains
non-significant ($p = 0.711$ HAC vs.\ $p = 0.695$ standard).

\textbf{Temporal stability and structural break.}
Pre-GFC Normal ($n = 3{,}140$): $p = 6.66\times10^{-16}$ ($\Delta R^2 = 2.06\%$).
Post-GFC Normal ($n = 1{,}557$): $p = 0.73$ ($\Delta R^2 < 0.01\%$).
(The sum $3{,}140 + 1{,}557 = 4{,}697$ is 26 days fewer than Table~\ref{tab:regimes}'s
Normal total of 4{,}723, due to standard regime-boundary exclusion at lag~1
removing the first observation of each Normal-regime episode.)
A GFC-motivated Chow test at January 2008 confirms this break: $F(3,n{-}6) = 9.68$,
$p = 2.29\times10^{-6}$; $\hat{\beta}_{\text{HML}}$ shifts from $-0.189$
(pre-GFC) to $+0.010$ (post-GFC, Wald $z = 5.05$, $p = 9.2\times10^{-7}$).
\emph{Caveat:} The January 2008 break date is narrative-motivated, not data-selected.
The Quandt-Andrews~\cite{andrews1993tests} sup-$F = 22.1$ exceeds the 1\% CV
($14.6$) for $q=3$, confirming the presence of \emph{at least one} structural break,
but the sup-$F$-maximising date is June 1998 (LTCM/Russia crisis), not January 2008.
This discrepancy suggests possible multi-break structure~\cite{bai1998estimating};
the Chow $p$-value at January 2008 is valid as a test of that specific date
but should not be interpreted as identifying the sole or primary breakpoint.
CUSUM ($p = 0.96$) is consistent with a sharp break.
A TOST equivalence test~\cite{schuirmann1987comparison} with margin $|\beta| < 0.05$
confirms the post-2008 coefficient is statistically equivalent to zero ($p_{\text{TOST}} < 0.001$),
strengthening the interpretation that the link genuinely vanished rather than merely
became undetectable post-2008.
Three non-exclusive hypotheses may explain the post-GFC attenuation.
First, \emph{regulatory restructuring}: the Volcker Rule and Dodd-Frank
constraints reduced proprietary trading in factor strategies, attenuating
the cross-factor spillovers that generated the HML$\to$SMB lag.
Second, \emph{factor crowding and correlation shifts}: post-crisis
institutional crowding into value and size strategies increased their
contemporaneous co-movement (unconditional HML--SMB correlation:
$\rho_{\text{pre}} = 0.01$ vs.\ $\rho_{\text{post}} = 0.23$),
potentially absorbing the one-day predictive lead into same-day effects.
Third, \emph{statistical artifact}: the full-sample Normal regime may
mix structurally different pre- and post-2008 episodes, and the
breakpoint detection (sup-$F$ at June 1998) hints at multi-break
structure that the single-break Chow test cannot fully resolve.
We cannot distinguish among these hypotheses with the current data.

\textbf{Incremental $R^2$.} Table~\ref{tab:r2} quantifies the economic magnitude
relative to an SMB AR baseline. The baseline AR(15) $R^2$ ranges from 0.59\%
(Elevated) to 1.01\% (Crisis), reflecting the low daily autocorrelation of
factor returns. Adding HML lags (1--15) to this baseline increases $R^2$ by
1.04\% in Normal, 1.32\% in Elevated, and 0.56\% in Crisis. Normal and Elevated pass conventional
significance levels ($p = 5.98 \times 10^{-5}$ and $p = 0.002$), while Crisis
does not ($p = 0.988$). In the reverse direction, SMB lags add 1.35\% in Crisis
with $p = 0.562$, so we do not claim directional confirmation there.

\begin{table}[t]
\centering
\caption{Incremental $R^2$ from adding HML lags (1--15) to SMB autoregression by regime.
$\Delta R^2$ = $R^2_{\text{unrestricted}} - R^2_{\text{restricted}}$.}
\label{tab:r2}
\small
\begin{tabular}{lcccc}
\toprule
Regime & $n$ & $R^2_{\text{AR}}$ & $\Delta R^2$ & $p$-value \\
\midrule
Normal & 4,361 & 0.61\% & 1.04\% & $5.98 \times 10^{-5}$ \\
Elevated & 2,690 & 0.59\% & 1.32\% & $0.002$ \\
Crisis & 981 & 1.01\% & 0.56\% & $0.988$ \\
\bottomrule
\end{tabular}
\end{table}

\begin{figure}[t]
\centering
\includegraphics[width=\columnwidth]{figures/granger_heatmap.pdf}
\caption{All-pairs Granger causality by regime ($-\log_{10}(p)$, lag=5).
Bordered cells with asterisks survive Bonferroni correction.
The strongest regime-differential pair is MKT-RF$\to$SMB; HML$\to$SMB appears as
a secondary differential pattern.}
\Description{Three heatmaps of pairwise Granger significance by regime; darker cells indicate lower p-values, with Bonferroni-surviving cells marked.}
\label{fig:heatmap}
\end{figure}

\textbf{All-pairs perspective.} Figure~\ref{fig:heatmap} shows Granger causality
$p$-values for all directed factor pairs across three regimes (lag~5, for visual
density across all 30 pairs; Table~\ref{tab:main} uses BIC-selected lag~1).
Several pairs show regime-dependent variation; MKT-RF$\to$SMB has the largest differential.
HML$\to$SMB ranks fourth in this differential metric, with strongest evidence in
Normal ($p = 1.0 \times 10^{-6}$ at lag 5), suggestive Elevated evidence
($p = 0.013$), and no Crisis evidence ($p = 0.880$). We therefore treat
HML$\to$SMB as one regime-sensitive relation among several, not the dominant pair.

\textbf{Lag specification and persistent structure.} We tested lags 1--15 for
HML$\to$SMB across regimes (Figure~\ref{fig:lag_sensitivity}). In Normal, the
effect is strongest at lag 1 and remains jointly significant when including up
to lag 15 (joint $F$-test of lags 1--15: $p = 5.98 \times 10^{-5}$,
Table~\ref{tab:r2}). Elevated shows weak-to-moderate joint evidence
(lags 1--15 joint $p = 0.002$), while Crisis is consistently
null (lag 1 $p = 0.695$, joint lag 1--15 $p = 0.988$). We present Elevated/Crisis lag
patterns as exploratory because they do not meet the Bonferroni-significant threshold.

\subsection{Economic Interpretation and Supporting Evidence}
\label{sec:interpretation}
\label{sec:overlap_evidence}
\label{sec:oos}
We hypothesize a deleveraging cascade channel: forced unwinding of value positions
simultaneously deleverages associated size exposures, creating a predictive lag from
HML to SMB (Appendix~\ref{app:mechanism}). FF25 portfolio overlap analysis (sensitivity
fit, seed~42) finds Granger significance concentrating in small-cap portfolios
($\rho_s = 0.35$, permutation $p = 0.046$; Appendix~\ref{app:overlap}), consistent
with this channel. Event-level validation across six stress episodes shows the expected
pattern in 2 of 6 events, with 4 of 6 directionally correct---heterogeneous and low-power
(Appendix~\ref{app:events}). Early warning detection is limited under the primary fit
(0\% for EU~2011 and COVID-2020; Appendix~\ref{app:earlywarning}).
Alternative explanations (common drivers, information propagation) cannot be ruled out.

\section{Discussion}

\subsection{Frozen-Parameter Out-of-Sample Test (Primary Validation)}
\label{sec:frozen_oos}

Two-stage approaches risk circularity: the HMM uses the same returns for
regime assignment that Granger tests subsequently analyze.
We resolve this by fitting the Student-$t$ HMM exclusively on 1990--2012
($n = 5{,}797$ days), freezing all parameters
($\boldsymbol{\mu}_k$, $\boldsymbol{\Sigma}_k$, $\nu_k$, $\mathbf{A}$),
and classifying 2013--2024 ($n = 3{,}020$ days) without refitting.
Granger tests are then run on the held-out labels using the same
\texttt{select\_lag\_bic} + \texttt{run\_granger\_at\_lag} pipeline as Table~\ref{tab:main},
ensuring full methodological consistency.

The frozen HMM assigns 24.0\% to Normal, 30.7\% to Elevated, and 45.3\%
to Crisis in the held-out period (vs.\ 21.5\%, 13.7\%, 64.8\% in training).
Notably, the Elevated-regime share more than doubles from 13.7\% (training)
to 30.7\% (test), consistent with the OOS Elevated significance.
After lag-1 boundary exclusion (same rule as Section~\ref{sec:granger_method}),
the OOS retains 661$+$836$+$1{,}299$=2{,}796$ of 3,020 days (7.4\% excluded).
The OOS exclusion rate (7.4\%) is higher than the full-sample lag-1 Granger exclusion
(0.67\%; 59/8,817), reflecting more frequent regime transitions when the
frozen 1990--2012 classifier is applied to the unseen 2013--2024 period.

\begin{table}[t]
\centering
\setlength{\tabcolsep}{5pt}
\caption{Frozen OOS Granger Causality (HML$\to$SMB), 2013--2024.
HMM trained on 1990--2012 only; relabeling uses train-period ordering;
lag $= 1$ pre-fixed from in-sample BIC; filtered probabilities (no future info).
HAC-$p$ uses Newey--West bandwidth~1 (matching Granger lag); Elevated HAC $p = 0.041$.}
\label{tab:frozen_granger}
\small
\begin{tabular}{@{}lcccc@{}}
\toprule
Regime & $n$ & $F$-$p$ & HAC-$p$ & $\Delta R^2$ \\
\midrule
Normal    & 661 & 0.149 & 0.153 & 0.32\% \\
Elevated  & 836 & \textbf{0.014} & \textbf{0.041} & \textbf{0.73\%} \\
Crisis    & 1,299 & 0.281 & 0.290 & 0.09\% \\
\bottomrule
\end{tabular}
\end{table}

\textbf{Result} (Table~\ref{tab:frozen_granger}). HML$\to$SMB shows suggestive evidence in the Elevated regime
($p = 0.014$; HAC $p = 0.041$; $\Delta R^2 = 0.73\%$ at lag~1),
while Normal ($p = 0.149$) and Crisis ($p = 0.281$) are null.
\emph{Convention note:} The $n = 836$ in Table~\ref{tab:frozen_granger} uses the
decimal-unit pipeline; the permutation test uses percentage-unit ($n = 953$), yielding
slightly different regime assignments. Both conventions yield $F$-$p < 0.05$.
This result does not survive Bonferroni correction across 30 pairs
(threshold $\alpha/30 = 0.00033$, family-wise $\alpha = 0.01$); its evidential weight
rests on the economic motivation for the HML--SMB pair
(Section~\ref{sec:pairselection}).
Notably, in-sample significance is expected given the degrees of freedom consumed
by regime identification, lag selection, and pair screening on the same data;
the OOS result is therefore the more informative---if more modest---test of
generalization.
A full 50-seed sensitivity
analysis finds that \textbf{2 of 3 local optima} show Elevated significance.
The 50 seeds converge to three local optima (train LL$\approx$138,002: 28 seeds,
$F$-$p < 0.05$, HAC-$p < 0.05$; LL$\approx$137,942: 13 seeds,
$F$-$p < 0.05$, HAC-$p < 0.05$; LL$\approx$137,712: 9 seeds, all null);
significance arises in the two higher-LL optima, each representing a distinct
HMM solution. Because seeds within an optimum converge to identical parameter estimates,
effective robustness is across 2 distinct optima rather than 41 independent initializations.
These results use forward-only (filtered)
HMM probabilities and train-period relabeling, eliminating all forms
of look-ahead and circularity.
OOS is treated as an economically motivated confirmatory test of the pre-identified pair
(Section~\ref{sec:pairselection}); the Bonferroni-corrected result is discussed therein.
\textbf{Power analysis.} At $n = 836$, the OOS test has 98.7\% power to detect an
in-sample-magnitude effect ($\Delta R^2 = 2.06\%$, Cohen $f^2 = 0.021$) and 69.7\%
power at the observed OOS magnitude ($\Delta R^2 = 0.73\%$, $f^2 = 0.0074$). The
effect size needed for 80\% power is $\Delta R^2 = 0.93\%$. The smaller OOS effect
therefore reflects genuine attenuation---not merely insufficient sample size---consistent
with the structural-break interpretation (pre-GFC link weakened, not disappeared entirely).
The reverse direction SMB$\to$HML in Normal is nominally significant ($F$-$p = 0.047$;
HAC-$p = 0.063$) but does not survive Bonferroni correction across 30 pairs and was
not pre-specified; we report it for completeness only.

\begin{table}[t]
\centering
\caption{Frozen OOS HML$\to$SMB across three local optima.
Each optimum represents a distinct HMM solution from the 50-seed sweep
(Section~\ref{sec:hmm}). Optima~A and B yield significant Elevated-regime
evidence; Optimum~C (lowest train-period LL) does not.}
\label{tab:optima_oos}
\small
\begin{tabular}{@{}lccccc@{}}
\toprule
Optimum & Seeds & Train LL & $n_{\text{Elev}}$ & $F$-$p$ & HAC-$p$ \\
\midrule
A (seed 20) & 28 & 138{,}002 & 931 & \textbf{0.008} & \textbf{0.026} \\
B (seed 28) & 13 & 137{,}942 & 836 & \textbf{0.014} & \textbf{0.041} \\
C (seed 6)  &  9 & 137{,}712 & 1{,}099 & 0.466 & 0.526 \\
\bottomrule
\end{tabular}
\end{table}

Table~\ref{tab:optima_oos} shows that the OOS Elevated result holds in the two higher-LL optima but not in the lowest-LL optimum. Because seeds within each optimum converge to identical parameters, effective robustness is across 2 distinct HMM solutions rather than 41 independent initializations.

\textbf{Regime interpretation.} The in-sample Normal-regime dominance
($p = 8.75 \times 10^{-9}$) vs.\ OOS Elevated significance reflects
regime redistribution under the frozen classifier: in training (1990--2012),
Crisis accounts for 64.8\% of days while Elevated is only 13.7\%;
in the held-out period (2013--2024), Crisis shrinks to 45.3\% while
Elevated grows to 30.7\%. The OOS period thus contains substantially more
Elevated-regime days than the training period, and the frozen classifier
concentrates the HML$\to$SMB signal in this regime.
The directional finding---HML predicts SMB during stress-adjacent
(Elevated) conditions---is consistent across the OOS design.

\textbf{Event-level detail.} Restricting to specific stress episodes
within the held-out period:

\begin{table}[t]
\centering
\caption{Held-Out Stress Events (2013--2024). HMM trained on 1990--2012 only.
Windows: 2015 China Aug--Oct~2015 (37 days), 2020 COVID Feb--Jun~2020 (94 Elevated-classified days; wider window than the 83 Crisis days in Appendix~\ref{app:events}), 2022 Rate Hikes Jan--Dec~2022.}
\label{tab:frozen_events}
\scriptsize
\begin{tabular}{@{}lcccc@{}}
\toprule
Event & Days & Crisis\% & $p$(H$\to$S) & $p$(S$\to$H) \\
\midrule
2015 China & 37 & 54\% & 0.165 & 0.123 \\
2020 COVID & 94 (Elevated) & 0\% & 0.032 & 0.107 \\
2022 Rate Hikes & 124 & 100\% & 0.472 & 0.186 \\
\bottomrule
\end{tabular}
\end{table}

Event-level behavior is heterogeneous (COVID $p = 0.032$; 2022 $p = 0.472$),
highlighting that the per-regime pooled result captures the aggregate
stress-period signal more reliably than individual events.

\subsection{Markov-Switching VAR Comparison}
\label{sec:msvar}
A single-stage MS-VAR confirms that HML lags contain information about SMB dynamics ($\text{LR} = 114.95$, $p < 10^{-12}$), but BIC favors the restricted model ($\Delta\text{BIC} = +130$). See Appendix~\ref{app:msvar}.

\subsection{Trading Strategy Evaluation}
\label{sec:trading}
A simple Crisis-regime trading strategy produces negative returns ($-0.3\%$ annual, Sharpe $= -0.07$), confirming that statistical predictability does not imply trading profitability. See Appendix~\ref{app:trading}.

\subsection{Risk Monitoring Application}
\label{sec:risk_monitoring}
A hybrid VaR detector (sensitivity fit, seed~42) combining HMM regime probabilities with a realized-volatility override achieves 5.60\% violation rate (Christoffersen CC $p = 0.336$), improving upon the unconditional baseline (6.42\%, $p = 0.002$). The improvement reflects primarily the volatility-override mechanism rather than the Granger link per se; the frozen OOS Granger result for seed~42 is null (Elevated $p = 0.466$). We therefore frame the VaR application as a proof-of-concept for regime-conditional monitoring architecture---the Granger signal informs the regime taxonomy that enables adaptive VaR windows, even though the direct HML$\to$SMB predictor does not drive the violation-rate improvement. See Appendix~\ref{app:var} for full VaR comparison, Algorithm~1, and hybrid detector details.

\subsection{Robustness}
\label{sec:robustness}

A 50{,}000-shuffle permutation test (Section~\ref{sec:pairselection}) yields
$p = 0.022$ (SE $= 0.0007$, 95\% CI [0.021, 0.023], $n = 953$)---significant at 5\%.
HAC bandwidth sensitivity (1, 4, 7, 10, 15) yields $p < 0.05$ throughout
(range: 0.015--0.025), confirming robustness to bandwidth choice.
The rank-2/30 result (Section~\ref{sec:pairselection}) confirms the
HML$\to$SMB pair is not a post-hoc artefact of in-sample screening.

\begin{figure}[t]
\centering
\includegraphics[width=\columnwidth]{figures/rolling_granger.pdf}
\caption{Rolling 3-year unconditional Granger causality ($-\log_{10}(p)$, HML$\to$SMB,
lag=9). Background shading shows HMM regime assignments for reference. Predictive
strength varies episodically, with peaks during periods overlapping known market
stress events. This rolling analysis does not condition on regime labels;
regime-conditional results are in Table~\ref{tab:main}.}
\Description{Time-series plot of rolling Granger significance with shaded regime bands, showing time variation in HML-to-SMB predictability.}
\label{fig:rolling}
\end{figure}

\textbf{Temporal dynamics.} Figure~\ref{fig:rolling} presents rolling 3-year
Granger $p$-values for HML$\to$SMB, showing that the predictive link is not a
static property, with peaks during calendar stress periods (2008--2009, 2020),
though this unconditional rolling test does not distinguish regime-conditional effects.
The signal is strongest during the 2008--2009 and 2020 crises and near-absent during
calm periods.

\textbf{Lag specification.} BIC selects lag 1 for all regime$\times$direction
combinations (Table~\ref{tab:main}).
Figure~\ref{fig:lag_sensitivity} shows significance is robust across lags 1--15
in Normal, while Crisis remains non-significant across lags.

\begin{figure}[t]
\centering
\includegraphics[width=\columnwidth]{figures/lag_sensitivity.pdf}
\caption{Granger causality $p$-values for HML$\to$SMB across lags 1--15 by regime.
Normal regime (green) shows concentrated short-lag significance ($p < 10^{-8}$ at lag 1).
Crisis regime (red) remains non-significant across the tested lag range.}
\Description{Line chart of HML-to-SMB p-values by lag for each regime, highlighting strong Normal-regime significance and weak Crisis-regime evidence.}
\label{fig:lag_sensitivity}
\end{figure}

\textbf{Trivariate robustness (MKT-RF control).}
To rule out a common market factor driving both HML and SMB with different lags,
we condition on one lag of MKT-RF in addition to SMB's own lag.
Results (frozen OOS, 2013--2024): Normal F-$p = 0.034$, HAC-$p = 0.041$;
Elevated F-$p = 0.012$, HAC-$p = 0.035$;
Crisis F-$p = 0.239$ (null). The Elevated signal \emph{strengthens} slightly after
controlling for MKT-RF (bivariate F-$p = 0.014$; trivariate F-$p = 0.012$; post-hoc),
and MKT-RF itself is non-significant in Elevated and Normal
(MKT-RF$\to$SMB: Elevated F-$p = 0.78$; Normal F-$p = 0.43$)
and marginally so in Crisis (F-$p = 0.077$).
The OOS Normal bivariate result is clearly null (HAC $p = 0.153$, F-$p = 0.149$).
Adding MKT-RF nominally yields trivariate HAC $p = 0.041$, but this was not
pre-specified, does not survive Bonferroni correction, and should not be interpreted
as evidence for OOS Normal predictive structure.
A common market factor therefore does not subsume the HML$\to$SMB predictive link.

\textbf{Regime count sensitivity ($K$).}
To assess whether the Granger result is an artifact of choosing $K = 3$, we fit
Student-$t$ HMMs with $K = 2$ and $K = 4$ on 1990--2012 training data
(10 seeds each, selecting by log-likelihood) and apply frozen OOS classification
to 2013--2024. Table~\ref{tab:k_sensitivity} reports the results with permutation
$p$-values (5{,}000 shuffles) for the lowest-$p$ regime at each $K$.
With $K = 2$, the signal is entirely absent (permutation $p = 0.572$).
With $K = 4$, the Low-Elevated regime approaches marginal significance
(F-$p = 0.056$) but does not survive permutation testing.
The $K = 3$ parameterization uniquely captures a distinct intermediate-volatility state
where predictive structure concentrates (permutation $p = 0.025$);
coarser ($K = 2$) or finer ($K = 4$) partitions dilute or fragment this signal.

\begin{table}[t]
\centering
\caption{$K$-Sensitivity: Frozen OOS HML$\to$SMB by Number of Regimes.
Best regime = lowest F-$p$ across regimes within each $K$.
Permutation: 5{,}000 label shuffles (percentage-unit convention).}
\label{tab:k_sensitivity}
\scriptsize
\begin{tabular}{@{}lclcccc@{}}
\toprule
$K$ & $\Delta$BIC$_{K\text{vs.}3}$ & Best Regime & $n$ & F-$p$ & Perm.\ $p$ & Sig.? \\
\midrule
2 & $+1{,}680$ & High-Vol & 1{,}917 & 0.514 & 0.572 & No \\
\textbf{3} & \textbf{---} & \textbf{Elevated} & \textbf{953} & \textbf{0.003} & \textbf{0.025} & \textbf{Yes} \\
4 & $+312$ & Low-Elev. & 855 & 0.056 & --- & No \\
\bottomrule
\end{tabular}
\end{table}

\textbf{Secondary factor pair validation (MOM$\to$SMB).}
To assess whether regime-dependent predictive structure is specific to HML$\to$SMB or a
more general phenomenon, we apply the identical frozen OOS methodology to MOM$\to$SMB---the
top-ranked pair by $F$-statistic in the OOS Elevated window ($F = 20.3$, rank 1/30;
Section~\ref{sec:pairselection}).
Results confirm the same regime-dependent pattern:
MOM$\to$SMB is significant in Elevated ($F = 16.5$, $F$-$p < 0.0001$, HAC $p = 0.003$,
$\Delta R^2 = 1.94\%$, $n = 836$; permutation $p = 0.010$, 10{,}000 shuffles, SE $= 0.001$)
but null in Normal ($F$-$p = 0.81$) and Crisis ($F$-$p = 0.15$).
The reverse direction (SMB$\to$MOM) is null in all regimes ($p > 0.57$), confirming
directional asymmetry.
This secondary validation strengthens the main finding: the phenomenon of
regime-conditional Granger predictive structure into SMB is not unique to HML but also
appears for Momentum, consistent with a common mechanism---possibly liquidity-driven
deleveraging that propagates across factors during elevated-volatility periods.

Additional robustness analyses (subsample stability, alternative regime methods, practitioner rolling-window baselines, weekly aggregation, regime-transition dynamics, neural complexity diagnostics, transfer entropy triangulation, and filtered vs.\ smoothed regime comparison) are provided in Appendix~\ref{app:robustness}.

\subsection{Interpretation and Scope}
\label{sec:limitations}

\textbf{Causal interpretation.}
Our core finding documents \emph{predictive precedence}, not structural causality:
lagged HML improves SMB prediction conditional on SMB's own history, but a common
driver (e.g., funding liquidity) affecting both factors with different lags cannot
be ruled out. The deleveraging cascade interpretation
(Appendix~\ref{app:mechanism}) is a plausible hypothesis, not a verified mechanism.
Our FF~25 overlap analysis provides portfolio-level evidence
($\rho_s = 0.35$, permutation $p = 0.046$; Appendix~\ref{app:overlap}),
but holdings-level verification via 13F data or fund flows remains for future work.
Formal pre-registration was not performed for the HML$\to$SMB pair;
while the economic prior (liquidity-mediated contagion) and rank-2/30
position (Section~\ref{sec:pairselection}) mitigate selection concerns,
this potential circularity cannot be fully resolved.

\textbf{Regime specification and $K$-sensitivity.}
The OOS Elevated result requires $K = 3$ (null under $K = 2$; marginal at $K = 4$,
$p = 0.056$). Training-period BIC strongly favors $K = 3$ ($\Delta\text{BIC} = 1{,}680$
over $K = 2$), but the effect's sensitivity to $K$ indicates it concentrates in
a specific intermediate-volatility state that coarser or finer partitions dilute.
This is an inherent limitation: regime-conditional findings are by definition
conditional on regime specification. Future work using Bayesian model averaging
over $K$ or nonparametric regime detection could test whether the effect
survives $K$-marginalization.

\textbf{Regime robustness.}
The frozen OOS design (Section~\ref{sec:frozen_oos}) eliminates
regime-identification circularity, but the 1990--2012-trained classifier
distributes OOS days differently than training: Elevated grows from 13.7\%
to 30.7\%, while Crisis shrinks from 64.8\% to 45.3\%.
Table~\ref{tab:optima_oos} shows that significance holds in 2 of 3 local optima
the 9 null seeds share the lowest train-period LL.
Distributional comparison (KS, Levene, $t$-test) finds no conditional-mean shift
between training and OOS Elevated ($t$-test $p \geq 0.15$), supporting regime
label stability for Granger inference.
However, KS ($p \approx 10^{-28}$) and Levene ($p \approx 10^{-175}$) reject
full distributional equivalence, indicating different tail behavior and dispersion.
Because the Granger $F$-test relies on conditional-mean structure rather than full distributional
identity, the mean-stable / tail-different pattern supports the validity of regime-conditional
Granger inference while cautioning that higher-moment properties (e.g., kurtosis, tail dependence)
may differ between periods; HAC standard errors partially address this heteroskedasticity.
Boundary cleaning at lag~1 excludes 7.4\% of OOS observations; Crisis contains
1,071 days in the primary fit (null at $p = 0.695$).

\textbf{Practical scope.}
The Normal-regime effect is exclusively pre-GFC ($p = 0.73$ post-2008;
Section~\ref{sec:main_result}). The OOS Elevated evidence is consistent with
regime migration but also with regime redistribution and $K$-specification
dependence. Economic magnitudes are small ($\Delta R^2 < 1\%$),
trading profitability is negative (Sharpe $= -0.07$; Appendix~\ref{app:trading}),
and VaR improvement (6.42\%$\to$5.60\%) reflects the volatility-override
mechanism rather than the Granger link (Appendix~\ref{app:var}).

\section{Conclusion}

Under the BIC-optimal primary fit ($K = 3$, seed~28), HML$\to$SMB Granger causality is regime-dependent and era-shifting: Bonferroni-significant in the Normal regime pre-GFC ($p = 8.75 \times 10^{-9}$, $\Delta R^2 = 2.06\%$), with a structural break at January 2008 (post-2008: $p = 0.73$). A frozen OOS design finds modest Elevated-regime evidence (permutation $p = 0.022$ [primary], $0.063$ [sensitivity]; 2 of 3 HMM optima), but this does not survive 3-regime Bonferroni correction and requires $K = 3$. A secondary pair (MOM$\to$SMB, permutation $p = 0.010$) confirms the regime-conditional pattern is not pair-specific. Transfer entropy reveals a reverse nonlinear SMB$\to$HML channel ($z = 5.37$), indicating the full predictive architecture is directionally asymmetric: linear forward, nonlinear reverse.

These findings establish predictive precedence, not structural causality; common drivers remain viable alternatives. Economic magnitudes are modest ($\Delta R^2 < 1\%$) and do not translate to trading profitability (Sharpe $= -0.07$). A hybrid VaR detector (Appendix~\ref{app:var}) demonstrates that regime-conditional architectures can improve tail-risk monitoring (5.60\% violation rate, CC $p = 0.336$), though the improvement is driven by the volatility-override mechanism rather than the Granger signal itself.

\textbf{Implications.}
Factor-timing models that assume regime-invariant cross-factor
relationships may misspecify dynamics; the HML$\to$SMB link that was
Bonferroni-significant pre-GFC has been absent for 16 years.
Regime-conditional architectures can improve tail-risk monitoring even when
the underlying statistical signal is modest.
The directional complexity asymmetry (linear forward, nonlinear reverse)
suggests standard Granger networks underestimate information flow.
Future work should pursue pre-registered replication on international factor
data, Bayesian model averaging over $K$, soft-label Granger estimation,
and holdings-level verification of the deleveraging channel via 13F data.

\section*{Funding}

This research received no specific grant from any funding agency in the public,
commercial, or not-for-profit sectors.

\section*{Data and Code Availability}

Code, analysis scripts, and the fixed seed list (50 seeds) for this paper are available at an anonymous repository to be provided upon acceptance. The financial data are publicly available from Kenneth French's data library at \url{https://mba.tuck.dartmouth.edu/pages/faculty/ken.french/data_library.html}. \emph{Reproducibility note:} Student-$t$ HMM estimation via EM is sensitive to initialization (tolerance: $\Delta\text{LL} < 10^{-6}$; max iterations: 500; all 50 seeds converged); independent researchers using different seeds may obtain different local optima. We provide fixed seeds and tie-breaking rules to enable exact replication of the primary fit (seed~28, selected as first in sorted order among three seeds reaching the same full-sample log-likelihood). The ``BIC-optimal'' label refers to $K = 3$ selection; the ``leakage-safe'' (LL-maximizing) label refers to seed selection within the $K = 3$ model, which avoids using crisis-event alignment as a selection criterion. The OOS Elevated result is significant in 2 of 3 local optima from our 50-seed sweep. Permutation test seeds are fixed and documented in the repository.

\appendix

\section{Hypothesized Economic Mechanism}
\label{app:mechanism}

We offer the following interpretation as a \emph{hypothesis}, not a verified mechanism:

\textbf{Normal-regime lag-1 link (pre-GFC, confirmatory):} During calmer periods,
cross-factor inventory management and short-horizon rebalancing may transmit
value shocks into size returns with a one-day delay.

\textbf{Stress-channel hypothesis (exploratory):} FF25 overlap is consistent with
a deleveraging channel in stress periods
($\rho_s = 0.35$, permutation $p = 0.046$). Primary Crisis-regime factor-level
Granger tests remain null in the leakage-safe fit.

\textbf{Caveat: dual-mechanism problem.} The Normal-regime ``inventory management''
story differs from the Elevated/Crisis ``deleveraging cascade'' hypothesis.
That different regimes require different mechanisms suggests either (a)~multiple
channels operate at different volatility levels, or (b)~the statistical pattern
reflects non-economic regularities. This ambiguity should temper confidence in
any single economic interpretation.

\textbf{Alternative explanations we cannot rule out:}
\begin{itemize}
    \item A common driver (e.g., funding liquidity) affects both factors with different lags
    \item Information about distress propagates through markets, reaching Value-sensitive
    investors before Size-sensitive investors
    \item Statistical artifact of regime-conditional estimation or $K = 3$ specification
\end{itemize}

Greenwood and Thesmar~\cite{greenwood2011stock} document that stocks with
concentrated institutional ownership exhibit correlated fragility during
liquidation events---consistent with our hypothesized deleveraging channel.
We provide exploratory portfolio-level evidence for this overlap in
Section~\ref{sec:overlap_evidence}.

\section{Portfolio Overlap Evidence}
\label{app:overlap}

The deleveraging hypothesis predicts that the HML$\to$SMB Granger link should
be strongest for portfolios that load on \emph{both} factors simultaneously.
The Small\-/HighBM portfolio sits in both the SMB long leg (small stocks) and
the HML long leg (high book-to-market), making it the primary candidate
transmission channel.

We test this prediction using the Fama-French 25 Size$\times$BM portfolios
(daily, 1990--2024). For each of the 25 portfolios, we run a Granger test
of HML$\to$Portfolio at lag 1 (BIC-selected in Table~\ref{tab:main}) within
the Crisis regime.

\textbf{Spatial gradient test (exploratory).} We define an ``overlap score'' for each
portfolio as $(4 - \text{size quintile}) + \text{BM quintile}$, ranging from
0 (Big/LowBM, no overlap) to 8 (Small/HighBM, maximal overlap). The
confirmatory test is a Spearman rank correlation between overlap score and
$-\log_{10}(p)$ of the Granger test across 25 portfolios, with significance
assessed by spatial permutation (10{,}000 shuffles of grid assignments;
time series are never touched).

\textbf{Result:} $\rho_s = 0.35$, permutation $p = 0.046$
(Figure~\ref{fig:ff25}). Granger significance concentrates in small-cap
portfolios: all five Small quintile portfolios are nominally significant
($p < 0.01$), with three surviving Bonferroni correction ($\alpha/25 = 0.002$):
S/L ($p = 2.5 \times 10^{-6}$, $\Delta R^2 = 1.47\%$),
S/3 ($p = 0.0014$), and S/H ($p \approx 0$, $\Delta R^2 = 1.79\%$).
One Big-quintile portfolio is nominally significant ($p < 0.01$), but none
survive Bonferroni. This spatial pattern is inconsistent
with a uniform factor-level effect and supports the overlap channel.

\textbf{Overlap fraction.} At lag 15 (matching Table~\ref{tab:r2}),
\[
\frac{\Delta R^2(\text{HML}\to\text{S/H})}{\Delta R^2(\text{HML}\to\text{SMB})}=0.392
\]
Small\-/HighBM alone accounts for 39\% of the HML$\to$SMB incremental $R^2$.

\begin{figure}[t]
\centering
\includegraphics[width=\columnwidth]{figures/ff25_overlap_granger_heatmap.pdf}
\caption{FF 25 Portfolio Overlap Analysis. Panel A: $-\log_{10}(p)$ for\\
HML$\to$Portfolio Granger tests (Crisis, lag=1). Significance concentrates in
Small quintile; $^{**}$Bonferroni-significant. Panel B: Spatial gradient---overlap
score vs.\ Granger significance ($\rho_s = 0.35$, permutation $p = 0.046$).}
\Description{Two-panel heatmap showing FF25 portfolio Granger significance by size and BM quintile, with spatial gradient plot.}
\label{fig:ff25}
\end{figure}

\section{Early Warning Assessment}
\label{app:earlywarning}

The Student-$t$ HMM detects crisis regimes before market peaks (Table~\ref{tab:warning}).
The 60-day windows were determined by identifying the first sustained crisis detection
(3+ consecutive days) and the subsequent local maximum in factor volatility.

\begin{table}[t]
\centering
\caption{Early Warning Lead Time. First Detection = first day of 3+ consecutive
Crisis regime assignments.}
\label{tab:warning}
\scriptsize
\begin{tabular}{@{}lccc@{}}
\toprule
Event & First Det. & Vol. Peak & Lead \\
\midrule
Lehman 2008 & May 9, 2008 & Sep 15, 2008 & 129 days \\
EU Crisis 2011 & None & Aug 8, 2011 & None \\
COVID 2020 & None & Mar 23, 2020 & None \\
\bottomrule
\end{tabular}
\end{table}

Under the primary fit, early warning is absent for EU Crisis 2011 and COVID 2020
(Table~\ref{tab:warning}), consistent with the primary fit's non-alignment to these
calendar events (Table~\ref{tab:detection}).

\section{Event-Based Consistency Checks}
\label{app:events}

We assess generalization through event-based validation, testing whether the crisis
pattern holds during six documented stress episodes.

\begin{table}[t]
\centering
\caption{Event-Based Consistency Checks. Expected pattern: HML $\to$ SMB significant during
Crisis. Result codes: \checkmark = expected pattern holds ($p < 0.10$ for HML$\to$SMB,
$p > 0.10$ for reverse); Dir.\ = correct direction but insufficient power;
$\times$ = opposite pattern.}
\label{tab:events}
\scriptsize
\begin{tabular}{@{}lcccc@{}}
\toprule
Event & Days & $p$(H$\to$S) & $p$(S$\to$H) & Result \\
\midrule
2008 Financial & 189 & 0.118 & 0.648 & Dir. \\
2011 EU Debt & 85 & 0.090 & 0.321 & \checkmark \\
2015 China & 64 & 0.419 & 0.459 & Dir. \\
2018 Vol Shock & 29 & 0.236 & 0.179 & $\times$ \\
2020 COVID & 83 & 0.062 & 0.134 & \checkmark \\
2022 Rate Hikes & 188 & 0.626 & 0.112 & $\times$ \\
\bottomrule
\end{tabular}
\end{table}

\textbf{Summary (exploratory, low-power):} The expected pattern (HML$\to$SMB
significant, reverse not) holds clearly in 2 of 6 events (2011, 2020). Two additional
events (2008, 2015)
show the correct direction ($p_{\text{HML}\to\text{SMB}} < p_{\text{SMB}\to\text{HML}}$)
but lack statistical power. Under the null hypothesis of no directional preference,
the probability of 4+ correct-direction events out of 6 by chance is approximately
0.11 (binomial test), though this test has limited discriminative power at $n = 6$.

\textbf{Exceptions.} Two events show reversed dynamics: 2018 Vol Shock and
2022 Rate Hikes (SMB$\to$HML direction). Both are short-duration or
policy-driven episodes distinct from the liquidity-driven crises (2008, 2011, 2020)
where the HML$\to$SMB direction is supported. This distinction between crisis
types is speculative---we lack direct evidence for the underlying mechanisms.

\textbf{Elevated regime finding.} In the primary fit, Elevated shows nominal
HML$\to$SMB evidence at lag 1 ($p = 0.004$) but this does not survive the
paper's Bonferroni threshold ($\alpha/30 = 0.00033$), and reverse direction
remains null ($p = 0.897$). Annual splits are mixed (67\% SMB$\to$HML),
reinforcing that Elevated evidence is suggestive rather than confirmatory.

\section{Markov-Switching VAR Comparison}
\label{app:msvar}

Our two-stage procedure (fit HMM, then test Granger) invites comparison with
Markov-Switching VARs~\cite{krolzig1997markov}, which jointly estimate
regime-dependent dynamics in a single stage.

We fit two Markov-Switching regression models for SMB with switching variance
($K = 3$ regimes, 9 lags): a \emph{restricted} model with only lagged SMB as
predictors, and an \emph{unrestricted} model that adds 9 lagged HML terms.

\textbf{Likelihood ratio test:} The unrestricted model significantly improves fit
($\text{LR} = 114.95$, $df = 27$, $p = 8.0 \times 10^{-13}$), confirming that
HML lags contain information about SMB dynamics not captured by SMB's own history,
even within a single-stage regime-switching framework.

\textbf{BIC comparison:} Despite the significant likelihood improvement, BIC favors
the restricted model ($\Delta\text{BIC} = -130$), reflecting the penalty for 27
additional parameters across 3 regimes. This indicates that while HML lags are
statistically informative, the improvement does not justify the model complexity
under a parsimony criterion.

\textbf{Implication.} The MS-VAR confirms that the HML--SMB predictive link exists
within a principled single-stage framework ($p < 10^{-12}$), but its economic magnitude
is too small to justify a complex regime-switching prediction model---consistent with
our $\Delta R^2 \approx 2\%$ finding and the trading strategy failure.

\section{Trading Strategy Backtest}
\label{app:trading}

We evaluate whether the predictive relationship translates to profits via a simple
strategy: during Crisis regime, go long SMB when HML return over the past 9 days
is positive, short SMB when negative.

\begin{table}[t]
\centering
\caption{Trading Strategy Backtest (1995--2024). Strategy trades SMB based on
9-day lagged HML signal during Crisis regimes only.}
\label{tab:trading}
\begin{tabular}{lcc}
\toprule
Metric & Strategy & Buy-and-Hold SMB \\
\midrule
Annual Return & $-0.3\%$ & $+0.1\%$ \\
Sharpe Ratio & $-0.07$ & $+0.06$ \\
Max Drawdown & $-30.0\%$ & $-43.5\%$ \\
\bottomrule
\end{tabular}
\end{table}

\textbf{Why does Granger causality not imply trading profits?}

The negative returns highlight an important distinction: \emph{statistical predictability}
$\neq$ \emph{economic predictability}. Several factors explain the gap:

\begin{enumerate}
    \item \textbf{Effect size:} The Granger test detects whether HML \emph{improves}
    SMB prediction, not whether the improvement is large. The $\Delta R^2$ from
    adding HML lags is 0.56\% in the Crisis regime (Table~\ref{tab:r2}); this is
    economically small and not statistically significant.

    \item \textbf{Sign vs.\ magnitude:} Granger causality tests whether lagged HML
    coefficients are jointly non-zero, not whether they reliably predict SMB \emph{direction}.
    The relationship may involve complex dynamics (e.g., mean reversion) that a simple
    directional strategy cannot capture.

    \item \textbf{Regime detection lag:} The HMM assigns regimes using filtered
    (smoothed) probabilities. Real-time regime detection would be noisier, introducing
    additional error.

    \item \textbf{Transaction costs:} Not modeled, but would further erode returns.
\end{enumerate}

\textbf{Hedging vs.\ alpha.} The practical value, if any, lies in \emph{risk monitoring}
rather than directional trading. When the HMM signals Crisis regime and HML shows
large negative returns, a risk manager might:
\begin{itemize}
    \item Increase attention to SMB exposure
    \item Tighten stop-losses on small-cap positions
    \item Pre-arrange liquidity for potential SMB hedging
\end{itemize}

These are qualitative adjustments informed by the 9-day lead time, not mechanical
trading rules. We formalize this intuition in Section~\ref{sec:risk_monitoring}.

\section{Risk Monitoring and VaR Application}
\label{app:var}

While the HML--SMB link fails as a trading signal, we test whether it improves
\emph{tail risk forecasting}. We compare five Value-at-Risk (VaR) models for SMB
at the 5\% level, trained on 1990--2012 and tested out-of-sample on 2013--2024
($n \approx 3{,}000$ days; exact counts vary by model due to regime-boundary
exclusion and warm-up requirements):

\begin{enumerate}
    \item \textbf{Unconditional:} 60-day rolling historical VaR.
    \item \textbf{GARCH(1,1):} One-step volatility forecast with parametric VaR.
    \item \textbf{Regime-conditional:} VaR window adapts to regime
    (Crisis: 30 days, Elevated: 50, Normal: 75).
    \item \textbf{HML-informed:} Regime-conditional, plus a $2.3\times$ stress
    multiplier when 9-day cumulative HML $< -0.5$ during Crisis regime.
    Parameters calibrated on training data only.
    \item \textbf{Hybrid (HMM+Vol):} HML-informed model plus a volatility
    override that forces Stress Alert when 20-day realized volatility exceeds
    the 95th percentile of the training period.
\end{enumerate}

\begin{table}[t]
\centering
\setlength{\tabcolsep}{4pt}
\caption[Out-of-Sample VaR Backtest (2013--2024)]{Out-of-Sample VaR Backtest (2013--2024). Target violation rate: 5\%.
Christoffersen CC tests conditional coverage (unconditional + independence).}
\label{tab:var}
\small
\begin{tabular}{lccc}
\toprule
Model & Viol.\ Rate & Dev.\ from 5\% & CC $p$-value \\
\midrule
Unconditional & 6.42\% & $+1.42$ pp & 0.002 \\
GARCH(1,1) & 3.91\% & $-1.09$ pp & 0.014 \\
Regime-Conditional & 6.69\% & $+1.69$ pp & $< 0.001$ \\
HML-Informed & 5.77\% & $+0.77$ pp & 0.171 \\
Hybrid (HMM+Vol) & \textbf{5.60\%} & \textbf{+0.60 pp} & \textbf{0.336} \\
\bottomrule
\end{tabular}
\end{table}
\setlength{\tabcolsep}{6pt}

Both the HML-informed and hybrid models pass the Christoffersen conditional
coverage test ($p > 0.05$); all other specifications are rejected. The hybrid
model (Algorithm~\ref{alg:monitor}) adds a realized-volatility override that
activates when 20-day realized vol exceeds the 95th percentile of training
data. This corrects the frozen HMM's COVID-2020 detection failure
(0\%$\to$47.8\% detection; Tables~\ref{tab:hybrid_sensitivity},
\ref{tab:var_hybrid}) while improving VaR coverage to 5.60\%
($p_{\text{CC}} = 0.336$). However, the hybrid detector activates on 14.9\% of
test days with a 93.2\% false-alarm rate---on most alert days no VaR violation
occurs. Practitioners must weigh correct tail coverage against unnecessary capital
reserves on non-violation alert days; we analyze the base-rate mechanics in the
appendix (Section~\ref{sec:risk_monitoring}, ``False alarms'' paragraph).

\textbf{Tick loss and model selection.}
Despite superior coverage, the HML-informed model has higher mean tick loss
(1.05 vs.\ 0.90 baseline). This apparent contradiction arises because tick loss
penalizes conservatism: a model with correct coverage but wider VaR incurs higher
tick loss on the 95\% of non-violation days. The Diebold-Mariano test, applied to
tick loss, would rank the under-covering baseline above the correctly-covering
HML-informed model---a known limitation of tick loss for comparing models at
different coverage levels. We report both metrics for transparency: tick loss
measures forecast sharpness, while Christoffersen CC measures correctness.

\textbf{False alarms and model complexity.}
In the HML-informed model, the stress multiplier activates on 376 days
(12.8\% of the test period) with a 93.4\% false alarm rate. In the hybrid model,
alerts activate on 440 days (14.9\%) with a 93.2\% false alarm rate---i.e., on
most alert days, no VaR violation occurs.
This is partly a base-rate artifact: with a 5\% violation target and 12.8\%
alert rate, at most 39\% of alerts can be ``true'' ($5/12.8$), yielding a
61\% false alarm floor. The 93.4\% exceeds this floor but reflects the cost
of correct tail coverage: alert-day VaR averages $2.51\times$ wider than
non-alert VaR, achieving zero violations on alert days at the cost of
unnecessary capital reserves on non-violation alerts.
The alternative (GARCH) avoids regime labels entirely and achieves
3.91\% violation rate with no false-alarm mechanism, but its over-coverage means
capital is consistently over-reserved. In practice, the choice depends on whether
the institution prefers fewer but louder alerts (HML-informed) or consistently
conservative risk limits (GARCH). Both outperform unconditional VaR; neither
dominates on all criteria.
We note that the HML-informed model's success may partly reflect conservative
VaR widening during high-volatility regimes rather than the Granger link per se,
since the underlying effect is non-significant in 2013--2024 ($p = 0.162$).

\textbf{COVID-19 common-mode failure and hybrid detector.}
The frozen HMM classifies the entire COVID period as non-Crisis (0\% Crisis days),
so the HML-informed stress multiplier never activates---a complete regime detection
failure. We address this with a \emph{hybrid detector} (Algorithm~\ref{alg:monitor})
that supplements HMM filtered probabilities with a 20-day realized volatility
fallback: when realized vol exceeds the 95th percentile of the \emph{training period}
(1990--2012), the detector overrides non-Crisis HMM classifications to ``Stress Alert.''
The volatility threshold ($\sigma_{95} = 1.446$) is a fixed design choice locked
before test evaluation---no search over thresholds on the holdout period.

\textbf{Hybrid detector results.} The volatility override activates on 44 of 92
COVID trading days (47.8\% detection rate vs.\ 0\% for HMM alone;
Table~\ref{tab:hybrid_sensitivity}). VaR performance improves further:
violation rate 5.60\% (deviation from 5\% target: $+0.60$ pp), Christoffersen
CC $p = 0.336$---the strongest conditional coverage pass among all specifications
(Table~\ref{tab:var_hybrid}; Figure~\ref{fig:hybrid_detector}). Average violation magnitude also improves
($-0.275$ vs.\ $-0.305$ for HML-Informed), indicating less severe tail breaches
when they occur.

\begin{table}[t]
\centering
\caption{Hybrid Detector: COVID-19 Detection Sensitivity. Vol threshold =
percentile of 20-day realized volatility computed on 1990--2012 training data.
Primary result uses 95th percentile (bolded).}
\label{tab:hybrid_sensitivity}
\small
\begin{tabular}{lcccc}
\toprule
Percentile & Threshold & Override Days & Detection Rate \\
\midrule
85th & 0.924 & 85 & 92.4\% \\
90th & 1.098 & 76 & 82.6\% \\
\textbf{95th} & \textbf{1.446} & \textbf{44} & \textbf{47.8\%} \\
99th & 2.122 & 25 & 27.2\% \\
\bottomrule
\end{tabular}
\end{table}

\begin{table}[t]
\centering
\caption{Out-of-Sample VaR Backtest (2013--2024) with Hybrid Detector.
Target violation rate: 5\%.}
\label{tab:var_hybrid}
\scriptsize
\begin{tabular}{@{}lcccc@{}}
\toprule
Model & Viol.\ Rate & Dev. & CC $p$ \\
\midrule
Unconditional & 6.42\% & $+1.42$ pp & 0.002 \\
HML-Informed & 5.77\% & $+0.77$ pp & 0.171 \\
\textbf{Hybrid (HMM+Vol)} & \textbf{5.60\%} & \textbf{+0.60 pp} & \textbf{0.336} \\
\bottomrule
\end{tabular}
\end{table}

\begin{figure}[t]
\centering
\includegraphics[width=\columnwidth]{figures/hybrid_detector_performance.pdf}
\caption{Hybrid detector performance. Panel A: realized volatility with 95th
percentile threshold and override periods (COVID window highlighted). Panel B:
SMB returns with unconditional, HML-informed, and hybrid VaR. Panel C:
cumulative violations vs.\ 5\% target.}
\Description{Three-panel figure showing realized volatility with override threshold, SMB returns with VaR bounds, and cumulative violation comparison.}
\label{fig:hybrid_detector}
\end{figure}

\begin{algorithm}[t]
\caption{Real-Time Factor Risk Monitor}
\label{alg:monitor}
\begin{algorithmic}[1]
\REQUIRE Frozen HMM parameters $(\boldsymbol{\mu}_k, \boldsymbol{\Sigma}_k, \nu_k, \mathbf{A})$;
vol threshold $\sigma_{95}$; HML threshold $\tau$; stress multiplier $m$;
regime windows $(w_0, w_1, w_2)$
\FOR{each trading day $t$}
  \STATE Compute filtered regime: $\hat{z}_t \gets \arg\max_k P(z_t=k \mid \mathbf{x}_{1:t})$
  \STATE Compute 20-day realized vol: $\hat{\sigma}_t \gets \text{std}(\|\mathbf{x}\|_{t-19:t})$
  \IF{$\hat{z}_t \neq \text{Crisis}$ \AND $\hat{\sigma}_t > \sigma_{95}$}
    \STATE $\hat{z}_t \gets \text{Crisis}$ \COMMENT{Volatility override}
  \ENDIF
  \STATE Set VaR window: $w \gets w_{\hat{z}_t}$
  \STATE Compute base VaR: $\text{VaR}_t \gets Q_{0.05}(\text{SMB}_{t-w:t-1})$
  \IF{$\hat{z}_t = \text{Crisis}$ \AND $\sum_{\ell=1}^{9}\text{HML}_{t-\ell} < \tau$}
    \STATE $\text{VaR}_t \gets m \cdot \text{VaR}_t$ \COMMENT{HML stress multiplier}
  \ENDIF
  \RETURN $\hat{z}_t$, $\text{VaR}_t$
\ENDFOR
\end{algorithmic}
\end{algorithm}

\begin{table}[t]
\centering
\caption{Algorithm~\ref{alg:monitor} Parameter Values. All calibrated on
training data (1990--2012) only.}
\label{tab:algo_params}
\small
\begin{tabular}{lcc}
\toprule
Parameter & Symbol & Value \\
\midrule
Vol threshold (95th pctile) & $\sigma_{95}$ & 1.446 \\
HML cumulative threshold & $\tau$ & $-0.5$ \\
Stress multiplier & $m$ & 2.3 \\
Window (Calm) & $w_0$ & 75 days \\
Window (Normal) & $w_1$ & 50 days \\
Window (Crisis) & $w_2$ & 30 days \\
HML lookback & --- & 9 days \\
\bottomrule
\end{tabular}
\end{table}

\section{Extended Robustness Analyses}
\label{app:robustness}

\textbf{Subsample stability.} Splitting at 2008: Crisis HML$\to$SMB
pre-2008 ($n = 659$, $p = 0.335$) and post-2008 ($n = 322$, $p = 0.599$)
are both non-significant, so we avoid claims of subsample-confirmed Crisis effects.

\textbf{Alternative regime methods.} Threshold-based volatility regimes
(realized volatility $>$ 90th percentile = Crisis) also detect the pattern
in sensitivity checks, but we treat this as exploratory given the primary-fit
null in Crisis at lag 1.

\textbf{Practitioner baselines (exploratory).} Standard unconditional tools detect the HML$\to$SMB
asymmetry but miss the architectural distinction. Rolling 250-day Granger tests
find HML$\to$SMB significant in 36.2\% of windows vs.\ 9.6\% for the
reverse---with concentration in Normal (47.1\% vs.\ 13.7\% reverse),
lower rates in Elevated (21.3\%), and mixed Crisis windows (27.7\%).
FEVD spillover is regime-insensitive in direction,
failing to detect the dormant Elevated regime or the Normal/Crisis directionality
shift. Regime conditioning does not merely sharpen detection---it reveals
qualitatively different predictive architectures that rolling methods average away.

\textbf{Weekly aggregation.} Using weekly returns, the HML$\to$SMB crisis pattern
remains non-significant ($p = 0.114$, $n = 221$), consistent with weak
evidence outside Normal.

\textbf{Regime transitions (exploratory).} Analyzing 60-day windows around Normal$\to$Crisis
transitions, the HML$\to$SMB relationship strengthens upon crisis entry
(median $p$: $0.058 \to 0.021$), suggesting discrete intensification of the
predictive link.

\textbf{Model complexity diagnostic protocol.}
We apply four model classes---Linear (OLS), RF (100 trees), MLP (64-32),
LSTM (32 hidden, 9-step)---to characterize
the \emph{complexity boundary} of the predictive relationship.
Each predicts SMB from its own lags (restricted) vs.\ both SMB and
HML lags (unrestricted), with permutation testing
(200 shuffles for Linear/RF/MLP, 100 for LSTM; at 200 permutations,
$p$-value SE $\approx 0.007$ at $p = 0.01$) and expanding-window CV.

Table~\ref{tab:neural} and Figure~\ref{fig:complexity} present the results.
Linear improvements are modest (Normal 0.86\%, Elevated 0.92\%, Crisis 0.43\%).
Under the primary fit (seed~28, canonical regimes $N = 2{,}506 / E = 4{,}670 / C = 1{,}641$),
\emph{no nonlinear method improves prediction in any regime}: RF, MLP, and LSTM
permutation tests are non-significant ($p > 0.16$ throughout).\footnote{Under the
sensitivity fit (seed~42, alternative regime partition $N = 2{,}506 / T = 4{,}670 / C = 1{,}641$),
RF shows significant improvement in the Transition regime ($p = 0.010$) and Crisis
($p = 0.005$), and MLP is marginally significant in Transition ($p = 0.015$).
This suggests the ``predominantly linear'' characterization holds for the primary
partition but not universally; nonlinear structure may emerge under alternative
regime definitions. The LSTM was estimated under a separate HMM initialization;
qualitative conclusions are consistent across fits.}
This complexity characterization supports a conservative conclusion:
incremental predictive content in the HML$\to$SMB direction is captured by
linear terms, though TE reveals residual nonlinear structure in the reverse
direction (Table~\ref{tab:te}).

\begin{figure}[t]
\centering
\includegraphics[width=\columnwidth]{figures/complexity_spectrum.pdf}
\caption{Four-model diagnostic protocol: MSE improvement (\%) from including
HML lags, by regime and model class. Only linear models (blue) achieve
significant improvement in Normal and Elevated; all nonlinear methods (RF, MLP, LSTM)
fail to improve prediction.}
\Description{Bar chart comparing MSE improvement across four model classes and three regimes, showing only linear models achieve significant improvement.}
\label{fig:complexity}
\end{figure}

\begin{table}[t]
\centering
\caption[Four-model diagnostics]{Four-Model Diagnostic Protocol: HML$\to$SMB by Regime.
Linear column shows MSE improvement (\%); others show permutation $p$-values.}
\label{tab:neural}
\small
\begin{tabular}{lcccccc}
\toprule
Regime & $n$ & Linear & RF $p$ & MLP $p$ & LSTM $p$ \\
\midrule
Normal & 4,496 & 0.86\%$^{**}$ & 0.69 & 0.20 & 0.63$^\dagger$ \\
Elevated & 2,792 & 0.92\%$^{**}$ & 1.00 & 1.00 & 0.93$^\dagger$ \\
\textbf{Crisis} & \textbf{1,017} & \textbf{0.43\%} & 0.13 & 0.96 & 0.83$^\dagger$ \\
\bottomrule
\multicolumn{7}{l}{\footnotesize $^{**}p < 0.01$. $^\dagger$LSTM: separate HMM init., 100 permutations.} \\
\end{tabular}
\end{table}

\textbf{Transfer entropy and directional complexity asymmetry.}
As a model-free check, we compute transfer entropy~\cite{schreiber2000measuring,barnett2009granger}
(Frenzel--Pompe kNN, $k=5$, 200 permutations).
Normal shows significant bidirectional TE (HML$\to$SMB $z = 2.45$, $p = 0.007$;
SMB$\to$HML $z = 5.37$, $p < 10^{-6}$). Elevated is weakly bidirectional at
conventional levels (HML$\to$SMB $p = 0.008$; SMB$\to$HML $p = 0.049$), while
Crisis is non-significant in both directions. The reverse-direction Normal TE
signal appears despite non-significant reverse linear Granger
($p = 0.864$, Table~\ref{tab:main}).

\textbf{Reconciling TE vs.\ Granger.} The directional asymmetry---linear
HML$\to$SMB but nonlinear SMB$\to$HML---is consistent with different
information transmission mechanisms. The forward HML$\to$SMB link
operates through a linear lag (e.g., inventory rebalancing with a
fixed delay); the reverse SMB$\to$HML channel may reflect nonlinear
feedback through liquidity or volatility states (e.g., small-cap stress
triggers disproportionate value-factor repricing only beyond a threshold).
This reconciles the null reverse Granger ($p = 0.864$) with the
highly significant reverse TE ($z = 5.37$): the reverse channel
is real but nonlinear, invisible to the linear Granger test.
We emphasize that the ``predominantly linear'' characterization applies
strictly to the \emph{forward} HML$\to$SMB direction under the primary fit;
the full bidirectional architecture is mixed-complexity.

\begin{table}[t]
\centering
\caption{Transfer Entropy: Multi-Lag (1--9) by Regime.
$z$-scores relative to permutation null (200 shuffles).}
\label{tab:te}
\scriptsize
\begin{tabular}{@{}lccccc@{}}
\toprule
Regime & $n$ & HML$\to$SMB $z$ & $p$ & SMB$\to$HML $z$ & $p$ \\
\midrule
Normal & 4,496 & $2.45^{**}$ & 0.007 & $5.37^{***}$ & $< 10^{-6}$ \\
Elevated & 2,792 & $2.41^{**}$ & 0.008 & $1.65^{*}$ & 0.049 \\
Crisis & 1,017 & $1.01$ & 0.157 & $1.22$ & 0.111 \\
\bottomrule
\multicolumn{6}{l}{\footnotesize $^{***}p < 10^{-6}$, $^{**}p < 0.01$, $^{*}p < 0.05$.} \\
\multicolumn{6}{l}{\footnotesize $n$ = clean observations after boundary exclusion.} \\
\end{tabular}
\end{table}

\textbf{Filtered vs.\ smoothed regime probabilities.} Smoothed probabilities use
the full sample (forward and backward pass); filtered probabilities use only past
data ($P(z_t = k \mid x_1, \ldots, x_t)$), simulating real-time classification.
Overall agreement is 95.9\%. Under filtered regimes, the Crisis HML$\to$SMB
result remains non-significant ($p = 0.601$ vs.\ $p = 0.695$ smoothed),
with filtered Crisis containing 1,079 days (vs.\ 1,071 smoothed).
The finding is qualitatively consistent across classification methods.

\section{Local Optima Taxonomy and Boundary Exclusion Details}
\label{app:optima}

\textbf{Local optima taxonomy.}
HMM estimation via EM is sensitive to initialization. Our 50-seed multistart
reveals 7 distinct local optima clusters:
Cluster~1 (3 seeds, best BIC~$= 75{,}587$, 0\% GFC detection, Crisis $p = 0.88$);
Cluster~2 (15 seeds, BIC~$= 75{,}625$, 0\% GFC, Crisis $p < 0.001$);
Cluster~3 (8 seeds, BIC~$= 75{,}660$, 0\% GFC, Crisis $p = 0.001$);
Cluster~4 (8 seeds, BIC~$= 75{,}726$, 0\% GFC, Crisis $p < 0.001$);
Cluster~5 (7 seeds, BIC~$= 75{,}805$, 90\% GFC, Crisis $p = 0.018$);
Cluster~6 (3 seeds, BIC~$= 75{,}906$, 100\% GFC, Crisis $p = 0.17$);
Cluster~7 (6 seeds, BIC~$= 76{,}137$, 92\% GFC, Crisis $p = 0.28$).
Only clusters with $\geq$50\% GFC detection (5--7) are economically valid for
crisis-regime inference; among these, Cluster~5 has the best fit.
The highest log-likelihood cluster (seeds 28, 20, 6; BIC~$= 75{,}587$)
assigns 0\% of 2008 crisis days to the Crisis regime ($p = 0.88$),
while a lower-likelihood cluster (seeds 21, 9, 0, 29, 46, 36, 11;
BIC~$= 75{,}805$, $\Delta\text{BIC} = 218$) assigns 90\% of 2008 to Crisis
($p = 0.018$). Readers who prefer pure BIC selection should interpret the
Crisis-regime Granger results as seed-dependent; the Elevated-regime OOS results
(primary seed 28) are robust across all 7 clusters ($p < 0.05$ in 5/7).

\textbf{Boundary exclusion statistics.}
With lag-1 Granger tests, boundary cleaning excludes only the first observation
of each new regime episode. In-sample exclusion rates:
Normal 0.55\% (26/4,723), Elevated 0.89\% (27/3,023), Crisis 0.56\% (6/1,071),
overall 0.67\% (59/8,817).
Frozen OOS exclusion is 7.4\% (224/3,020), reflecting more frequent
regime transitions in the 2013--2024 test period.
Calendar-day gaps between consecutive same-regime observations confirm that
excluded observations are structurally minor: median gap 3~calendar days,
95th percentile 4~days (weekend and holiday breaks);
only 1.2--3.3\% of inter-regime gaps exceed 20 trading days ($\approx$1 month).
Maximum gaps (1,568 days Normal, 2,423 days Elevated, 4,957 days Crisis)
reflect rare intervals between separate regime episodes.
Exclusion rates are not systematically biased toward any particular regime.

\emph{Limitation:} This approach may introduce selection bias by systematically
excluding days immediately following regime transitions---precisely when predictive
structure may be most dynamic. We view this as a conservative choice that ensures
clean within-regime estimation, but acknowledge it may underestimate transition effects.

\bibliography{references}

\end{document}
